% A font expanded automatically is a font of its base's own: the copies
% microtype makes for its contexts (\pdfcopyfont) each load theirs at
% the same ratio, and pdfTeX numbers each, as the /F name of any font
% loaded after them shows.
\documentclass{article}
\pdfadjustspacing=2
\font\base=cmr10
\pdfcopyfont\one=\base
\pdfcopyfont\two=\base
\pdfcopyfont\three=\base
\pdffontexpand\one 20 20 5 autoexpand
\pdffontexpand\two 0 20 5 autoexpand
\pdffontexpand\three 5 5 5 autoexpand
\font\later=cmbx10
\begin{document}
\one This paragraph, set in the first copy, is long enough to break into
a few lines, and each is stretched or shrunk a little by the expansion.
\par
\two This paragraph, set in the second copy, is long enough to break into
a few lines, and each is stretched or shrunk a little by the expansion.
\par
\three This paragraph, set in the third copy, is long enough to break into
a few lines, and each is stretched or shrunk a little by the expansion.
\par
\later The font loaded after them: \pdffontname\later.
\end{document}
